\documentclass[11pt]{article}
\usepackage[margin=1in]{geometry}
\usepackage{enumitem}
\usepackage{longtable}
\usepackage{booktabs}
\usepackage{amsmath,amssymb}
\usepackage{hyperref}
\hypersetup{colorlinks=true, linkcolor=blue, citecolor=blue, urlcolor=blue}

\newcommand{\Q}{\mathbb{Q}}
\newcommand{\Z}{\mathbb{Z}}
\newcommand{\A}{\mathbb{A}}
\newcommand{\Qp}{\mathbb{Q}_p}
\newcommand{\Zp}{\mathbb{Z}_p}
\newcommand{\OK}{\mathcal{O}_K}
\newcommand{\OL}{\mathcal{O}_L}
\newcommand{\Gal}{\mathrm{Gal}}

\title{Number Theory 1 \\ \large Course Design Package (PhD Level)}
\author{Lecture Format --- 14--15 Weeks, Two Sessions per Week}
\date{}

\begin{document}
\maketitle
\tableofcontents

\section*{Course Summary}
\begin{itemize}[noitemsep]
  \item \textbf{Course name:} Number Theory 1
  \item \textbf{Level:} PhD
  \item \textbf{Structure:} Lecture
  \item \textbf{Duration:} 14--15 weeks, two sessions per week ($\approx$28--30 sessions)
  \item \textbf{Assessment:} Weekly homework only; no exams, no final
  \item \textbf{Assumed background:} Passed algebra and analysis qualifying exams; real and complex analysis including measure theory (Haar measure); Galois theory
\end{itemize}

\section*{Explicit Assumptions}
\begin{itemize}[noitemsep]
  \item ``Algebra qualifying exam'' is taken to include standard commutative algebra (Noetherian rings, localization, integral closure, tensor products), since this is not stated explicitly but is required for Module 1 and not re-derived here.
  \item The 14--15 week range is scheduled as 15 weeks / 30 sessions, with Module 4 (Tate's thesis) placed last and treated as compressible, consistent with its ``if time permits'' status in the topic list.
  \item Weekly homework is assumed to mean one problem set per week (not per session), aligned to that week's module content.
\end{itemize}

%=====================================================================
\section{Module 1: Dedekind Domains and Ramification}
\textbf{Weeks 1--4 (Sessions 1--8).} \textit{Goal:} establish the ring-theoretic foundations --- Dedekind domains, ideal factorization, and ramification in extensions --- that underlie everything else in the course.

%---------------------------------------------------------------------
\subsection{Dedekind Domains: Definitions and Characterizations}
\textbf{Core concepts}
\begin{itemize}[noitemsep]
  \item Noetherian, integrally closed, Krull dimension one (every nonzero prime is maximal)
  \item Equivalent characterizations: localization at every prime is a DVR; every nonzero fractional ideal is invertible
  \item Examples: $\OK$ for a number field $K$, coordinate rings of smooth affine curves, PIDs
  \item Non-example: non-maximal orders, e.g.\ $\Z[\sqrt{-3}]$
\end{itemize}
\textbf{Learning objectives}
\begin{itemize}[noitemsep]
  \item Explain the equivalent characterizations of a Dedekind domain and prove their equivalence
  \item Derive that $\OK$ is a Dedekind domain from integrality and finiteness properties
  \item Apply localization arguments to reduce global ring-theoretic questions to the DVR case
\end{itemize}
\textbf{Example questions}
\begin{enumerate}[noitemsep]
  \item[Q1 (foundational).] Show that a PID is a Dedekind domain. \emph{Rubric: PID $\Rightarrow$ Noetherian (ACC on ideals); UFD $\Rightarrow$ integrally closed; nonzero primes in a PID are generated by irreducibles, hence maximal.}
  \item[Q2 (advanced).] Let $A = \Z[\sqrt{-3}]$. Exhibit an element of $\mathrm{Frac}(A)$ integral over $A$ but not in $A$, and conclude $A$ is not Dedekind. \emph{Answer: $\tfrac{1+\sqrt{-3}}{2}$ satisfies $x^2-x+1=0$ and lies in $\mathcal{O}_{\Q(\sqrt{-3})}=\Z[\tfrac{1+\sqrt{-3}}{2}]$, but not in $A$.}
\end{enumerate}

%---------------------------------------------------------------------
\subsection{Fractional Ideals, Unique Factorization, and the Ideal Class Group}
\textbf{Core concepts}
\begin{itemize}[noitemsep]
  \item Fractional ideals and their group structure under multiplication
  \item Unique factorization of ideals into primes (existence and uniqueness)
  \item The ideal class group $\mathrm{Cl}(K) = $ fractional ideals modulo principal ideals
  \item Finiteness of $\mathrm{Cl}(K)$ (stated here; proved in Module 2 via Minkowski theory)
\end{itemize}
\textbf{Learning objectives}
\begin{itemize}[noitemsep]
  \item Prove unique factorization of ideals in a Dedekind domain
  \item Compute the ideal class group in small examples
  \item Explain the relationship between class number one and unique factorization of elements
\end{itemize}
\textbf{Example questions}
\begin{enumerate}[noitemsep]
  \item[Q1 (foundational).] Show every nonzero fractional ideal of a Dedekind domain is invertible. \emph{Rubric: prove local invertibility at each prime using the DVR structure, then glue.}
  \item[Q2 (advanced).] Compute $\mathrm{Cl}(\Q(\sqrt{-5}))$ by exhibiting a non-principal ideal. \emph{Rubric: factor $(2,1+\sqrt{-5})$ as a prime above 2, show non-principal via a norm argument, show its square is principal; conclude $\mathrm{Cl}(K)\cong \Z/2$.}
\end{enumerate}

%---------------------------------------------------------------------
\subsection{Extensions of Dedekind Domains: Decomposition of Primes}
\textbf{Core concepts}
\begin{itemize}[noitemsep]
  \item The integral closure of a Dedekind domain in a finite extension is again Dedekind
  \item Factorization $p\OL = \prod P_i^{e_i}$ in an extension $L/K$
  \item The ideal norm and its multiplicativity
\end{itemize}
\textbf{Learning objectives}
\begin{itemize}[noitemsep]
  \item Derive the factorization $p\OL = \prod P_i^{e_i}$ from unique factorization
  \item Compute prime factorizations in specific number field extensions
  \item Apply the norm map to relate ideal factorization to field degree
\end{itemize}
\textbf{Example questions}
\begin{enumerate}[noitemsep]
  \item[Q1 (foundational).] Factor $(2)$ in $\Z[i]$. \emph{Answer: $(2) = (1+i)^2$, ramified since $2 \mid \mathrm{disc}(\Q(i))$.}
  \item[Q2 (advanced).] Determine how $5$ and $11$ factor in $\Q(\zeta_5)$. \emph{Rubric: $5$ is totally ramified since $5\mid \mathrm{disc}(\Q(\zeta_5))$; the splitting of $11$ is governed by the order of $11$ mod $5$.}
\end{enumerate}

%---------------------------------------------------------------------
\subsection{Ramification Index, Residue Degree, and the Fundamental Identity}
\textbf{Core concepts}
\begin{itemize}[noitemsep]
  \item Definitions of ramification index $e_i$ and residue degree $f_i$
  \item The fundamental identity $\sum e_i f_i = [L:K]$; constant $e,f$ when $L/K$ is Galois
  \item Split, inert, and ramified primes, with examples
\end{itemize}
\textbf{Learning objectives}
\begin{itemize}[noitemsep]
  \item Derive $\sum e_i f_i = n$ from the structure of $\OL \otimes_{\OK} \OK/p$
  \item Classify primes as split/inert/ramified in quadratic and cyclotomic fields
  \item Apply the identity to compute decomposition types from factorization data
\end{itemize}
\textbf{Example questions}
\begin{enumerate}[noitemsep]
  \item[Q1 (foundational).] Verify $\sum e_if_i = n$ for the factorization of $(2)$ in $\Z[i]$. \emph{Answer: $e=2, f=1, n=2$: $2\cdot 1 = 2$. \checkmark}
  \item[Q2 (advanced).] Give a non-Galois cubic field where the primes above some $p$ have differing $e_i,f_i$. \emph{Rubric: e.g.\ $x^3-x-1$; in a non-Galois extension $e_i,f_i$ need not be constant across the primes above $p$, unlike the Galois case.}
\end{enumerate}

%---------------------------------------------------------------------
\subsection{Discriminant, Different, and Dedekind's Ramification Theorem}
\textbf{Core concepts}
\begin{itemize}[noitemsep]
  \item The different ideal $\mathfrak{d}_{L/K}$ and its relation to ramification
  \item The discriminant as the norm of the different
  \item Dedekind's theorem: $p$ ramifies in $L$ iff $p \mid \mathrm{disc}(L/K)$
  \item Tame versus wild ramification
\end{itemize}
\textbf{Learning objectives}
\begin{itemize}[noitemsep]
  \item Explain the relationship between the different, the discriminant, and ramification
  \item Derive Dedekind's ramification criterion from properties of the different
  \item Distinguish tame from wild ramification and identify which primes can be wildly ramified
\end{itemize}
\textbf{Historically hard.} \emph{Why:} notational overload (different vs.\ discriminant, ideal norm vs.\ field norm) compounds with a genuine abstraction jump into wild ramification, which has no analogue in the tame case. \emph{Mitigation:} work tame ramification fully with explicit quadratic-field examples first; use a side-by-side table contrasting different vs.\ discriminant before stating Dedekind's theorem in general.
\textbf{Example questions}
\begin{enumerate}[noitemsep]
  \item[Q1 (foundational).] Compute $\mathrm{disc}(\Q(\sqrt d))$ for squarefree $d$ and identify the ramified primes. \emph{Answer: $\mathrm{disc}=d$ or $4d$ depending on $d \bmod 4$; ramified primes are exactly those dividing it.}
  \item[Q2 (advanced).] Explain why wild ramification at $p$ requires $p \mid e$, and construct an example. \emph{Rubric: tame ramification needs $e$ coprime to the residue characteristic; $\Q_p(p^{1/p})/\Q_p$ via the Eisenstein polynomial $x^p-p$ is totally, wildly ramified.}
\end{enumerate}

%---------------------------------------------------------------------
\subsection{Galois Extensions: Decomposition and Inertia Groups, Frobenius}
\textbf{Core concepts}
\begin{itemize}[noitemsep]
  \item Decomposition group $D_P$ and inertia group $I_P$ for $P$ above $p$ in Galois $L/K$
  \item Exact sequence $1 \to I_P \to D_P \to \Gal(k(P)/k(p)) \to 1$
  \item The Frobenius element/conjugacy class at unramified primes
  \item Statement of Chebotarev density, for motivation
\end{itemize}
\textbf{Learning objectives}
\begin{itemize}[noitemsep]
  \item Derive the structure of $D_P$ and $I_P$ from the Galois action on primes above $p$
  \item Compute Frobenius elements in explicit Galois extensions
  \item Explain how Frobenius conjugacy classes encode splitting behavior
\end{itemize}
\textbf{Historically hard.} \emph{Why:} students must track three interacting group actions (the full Galois group, $D_P$, $I_P$) simultaneously, and the analogy with the residue field's Galois group is easy to lose. \emph{Mitigation:} use one running example (e.g.\ $\Q(\zeta_n)$) throughout, with explicit subgroup-lattice diagrams at each step.
\textbf{Example questions}
\begin{enumerate}[noitemsep]
  \item[Q1 (foundational).] Compute $D_P, I_P$ for a split, an inert, and a ramified prime in a quadratic field. \emph{Answer: split: $D=I=1$; inert: $D=\Z/2, I=1$; ramified: $D=I=\Z/2$.}
  \item[Q2 (advanced).] Compute the Frobenius at $p=2$ in $\Q(\zeta_7)/\Q$ and determine the splitting of $2$. \emph{Rubric: order of $2$ mod $7$ is $3$, so $f=3$; three unramified primes each with $f=3$.}
\end{enumerate}

%=====================================================================
\section{Module 2: Global and Local Fields}
\textbf{Weeks 5--9 (Sessions 9--18).} \textit{Goal:} develop local fields as the local counterpart to number fields, and connect local and global perspectives via places, Minkowski theory, and units.

%---------------------------------------------------------------------
\subsection{Number Fields and Global Fields: Rings of Integers Revisited}
\textbf{Core concepts}
\begin{itemize}[noitemsep]
  \item Global fields: number fields and function fields over finite fields, as a unifying perspective
  \item Places of a global field: archimedean and non-archimedean
  \item $\OK$'s ideal theory recast in the global-field/local-field language used going forward
\end{itemize}
\textbf{Learning objectives}
\begin{itemize}[noitemsep]
  \item Explain the analogy between number fields and function fields as global fields
  \item Enumerate the places of a given number field (real, complex, finite)
  \item Situate $\OK$'s ideal theory within the local/global framework
\end{itemize}
\textbf{Example questions}
\begin{enumerate}[noitemsep]
  \item[Q1 (foundational).] List the places of $\Q(i)$. \emph{Answer: one complex place ($r_1=0, r_2=1$); non-archimedean places correspond to primes of $\Z[i]$.}
  \item[Q2 (advanced).] Compare places of $\mathbb{F}_q(t)$ with places of $\Q$, and identify the ``infinite place'' of $\mathbb{F}_q(t)$. \emph{Rubric: places of $\mathbb{F}_q(t)$ correspond to monic irreducibles plus the degree valuation at infinity; note the disanalogy that $\mathbb{F}_q(t)$ has no true archimedean place.}
\end{enumerate}

%---------------------------------------------------------------------
\subsection{Absolute Values and Completions; $p$-adic Numbers}
\textbf{Core concepts}
\begin{itemize}[noitemsep]
  \item Absolute values on a field; equivalence; Ostrowski's theorem for $\Q$
  \item Completion with respect to an absolute value; construction of $\Qp$
  \item $\Zp$ as an inverse limit; $p$-adic expansions
\end{itemize}
\textbf{Learning objectives}
\begin{itemize}[noitemsep]
  \item Prove Ostrowski's theorem classifying absolute values on $\Q$
  \item Construct $\Qp$ both via completion and via inverse limits, and show these agree
  \item Apply $p$-adic expansions to solve concrete equations
\end{itemize}
\textbf{Example questions}
\begin{enumerate}[noitemsep]
  \item[Q1 (foundational).] Verify $|\cdot|_p$ satisfies the ultrametric inequality. \emph{Rubric: follows from additivity properties of the $p$-adic valuation $v_p$.}
  \item[Q2 (advanced).] Does $\sqrt 2$ exist in $\Q_7$? If so, find its first three $p$-adic digits. \emph{Answer: $3^2 \equiv 2 \bmod 7$, so by Hensel's lemma $\sqrt 2 \in \Q_7$; digits obtained by Hensel lifting.}
\end{enumerate}

%---------------------------------------------------------------------
\subsection{Structure of Local Field Extensions; Hensel's Lemma}
\textbf{Core concepts}
\begin{itemize}[noitemsep]
  \item Hensel's Lemma and consequences (lifting roots, structure of $\Zp^*$)
  \item Uniqueness of the extension of an absolute value to a finite extension
  \item Unramified extensions $\leftrightarrow$ residue field extensions; totally ramified extensions $\leftrightarrow$ Eisenstein polynomials
  \item Structure theorem: every finite extension decomposes as unramified-then-totally-ramified
\end{itemize}
\textbf{Learning objectives}
\begin{itemize}[noitemsep]
  \item Apply Hensel's Lemma to lift factorizations and determine unit structure
  \item Classify finite extensions of $\Qp$ into unramified and totally ramified pieces
  \item Construct explicit local extensions via Eisenstein polynomials
\end{itemize}
\textbf{Historically hard.} \emph{Why:} the ``lifting'' intuition (a $p$-adic Newton's method) is unfamiliar, and holding the unramified/totally-ramified dichotomy in mind as two distinct extension mechanisms is easy to conflate. \emph{Mitigation:} present Hensel's Lemma explicitly as $p$-adic Newton's method with a numerical example first, then build the dichotomy from two contrasting explicit extensions.
\textbf{Example questions}
\begin{enumerate}[noitemsep]
  \item[Q1 (foundational).] Use Hensel's Lemma to show $\Zp^*$ contains all $(p{-}1)$-th roots of unity, $p$ odd. \emph{Rubric: apply Hensel to $x^{p-1}-1$ using simple roots mod $p$ from Fermat's little theorem.}
  \item[Q2 (advanced).] Construct the unramified extension of $\Qp$ of degree $n$, and contrast with $\Qp(p^{1/n})$. \emph{Rubric: unramified: adjoin roots of unity generating the degree-$n$ residue field extension; ramified: Eisenstein polynomial $x^n-p$.}
\end{enumerate}

%---------------------------------------------------------------------
\subsection{Local-Global Compatibility: Places, Completions, and the Product Formula}
\textbf{Core concepts}
\begin{itemize}[noitemsep]
  \item Normalized absolute values at each place; the product formula $\prod_v |x|_v = 1$
  \item Relating global ideal factorization to local completions $K_v$
  \item Weak approximation
\end{itemize}
\textbf{Learning objectives}
\begin{itemize}[noitemsep]
  \item Prove the product formula for number fields
  \item Explain how localizing $\OK$ at a prime recovers the local field $K_v$
  \item Apply weak approximation to construct elements with prescribed local behavior
\end{itemize}
\textbf{Example questions}
\begin{enumerate}[noitemsep]
  \item[Q1 (foundational).] Verify $\prod_v |2|_v = 1$ for $K = \Q$. \emph{Answer: $|2|_2 = 1/2$, $|2|_\infty = 2$, all other $|2|_v = 1$; product $=1$.}
  \item[Q2 (advanced).] Use the product formula to constrain elements of $\OK$ that are small at every archimedean place and integral at every finite place. \emph{Rubric: the product formula forces a global product-of-one constraint; relate to Kronecker's theorem characterizing roots of unity.}
\end{enumerate}

%---------------------------------------------------------------------
\subsection{Minkowski Theory: Geometry of Numbers, Finiteness of the Class Number}
\textbf{Core concepts}
\begin{itemize}[noitemsep]
  \item The Minkowski embedding $K \hookrightarrow \mathbb{R}^{r_1}\times \mathbb{C}^{r_2}$ and the lattice $\OK$
  \item Minkowski's convex body theorem
  \item The Minkowski bound; finiteness of $\mathrm{Cl}(K)$, completing Module 1
\end{itemize}
\textbf{Learning objectives}
\begin{itemize}[noitemsep]
  \item Derive the Minkowski bound and use it to compute class numbers in examples
  \item Prove finiteness of the ideal class group via the geometry of numbers
  \item Apply the Minkowski embedding to bound discriminants
\end{itemize}
\textbf{Historically hard.} \emph{Why:} this is a genuine change of category, from ideal-theoretic to geometric/volume-based reasoning, and the convex body theorem's soft counting argument is unlike the explicit proofs used so far. \emph{Mitigation:} draw the 2-dimensional lattice picture for a real quadratic field before the general proof; give a step-by-step worked template for the Minkowski bound.
\textbf{Example questions}
\begin{enumerate}[noitemsep]
  \item[Q1 (foundational).] Compute the Minkowski bound for $\Q(\sqrt{-5})$ and determine the class number. \emph{Answer: bound $\approx 2.85$, so only $p=2$ needs checking, giving class number $2$.}
  \item[Q2 (advanced).] Prove Minkowski's convex body theorem. \emph{Rubric: pigeonhole argument on translates of half the body; volume exceeding the covolume forces overlap.}
\end{enumerate}

%---------------------------------------------------------------------
\subsection{Dirichlet's Unit Theorem}
\textbf{Core concepts}
\begin{itemize}[noitemsep]
  \item The unit group $\OK^*$: torsion (roots of unity) and free part
  \item Statement: rank $r_1 + r_2 - 1$
  \item The logarithmic embedding and the geometry-of-numbers proof strategy
\end{itemize}
\textbf{Learning objectives}
\begin{itemize}[noitemsep]
  \item State and explain the structure of $\OK^*$ via Dirichlet's unit theorem
  \item Derive the rank formula $r_1+r_2-1$ using the logarithmic embedding
  \item Compute fundamental units in small examples
\end{itemize}
\textbf{Historically hard.} \emph{Why:} the logarithmic embedding reuses the Minkowski-theory machinery in a new embedding, and bookkeeping the rank $r_1+r_2-1$ against the torsion kernel is a common source of off-by-one errors. \emph{Mitigation:} work the rank computation for a real and an imaginary quadratic field side by side before the general proof.
\textbf{Example questions}
\begin{enumerate}[noitemsep]
  \item[Q1 (foundational).] Compute the unit rank of $\Q(\sqrt2)$ and $\Q(\sqrt{-2})$, and find a fundamental unit for $\Q(\sqrt2)$. \emph{Answer: ranks $1$ and $0$ respectively; fundamental unit $1+\sqrt2$.}
  \item[Q2 (advanced).] Derive $r_1+r_2-1$ from the logarithmic embedding $\OK^*\to \mathbb{R}^{r_1+r_2}$. \emph{Rubric: image lies in a hyperplane by the product formula; discreteness plus boundedness gives a full-rank lattice in that hyperplane.}
\end{enumerate}

%=====================================================================
\section{Module 3: Adèles and Idèles}
\textbf{Weeks 10--12 (Sessions 19--24).} \textit{Goal:} assemble all completions of a global field into the adèle ring and idèle group, developing the topological and measure-theoretic structure needed for Tate's thesis.

%---------------------------------------------------------------------
\subsection{Restricted Direct Products and the Adèle Ring}
\textbf{Core concepts}
\begin{itemize}[noitemsep]
  \item The restricted direct product construction $\prod'(K_v, \mathcal{O}_v)$
  \item The adèle ring $\A_K$ as a restricted product over all places
  \item $K$ embeds diagonally in $\A_K$; $K$ is discrete (cocompactness stated, proved in \S3.2)
\end{itemize}
\textbf{Learning objectives}
\begin{itemize}[noitemsep]
  \item Construct $\A_K$ as a restricted direct product and verify it is a topological ring
  \item Explain why restricting to ``almost all $\mathcal{O}_v$'' is necessary for local compactness
  \item Identify $K$ as a discrete subring of $\A_K$
\end{itemize}
\textbf{Historically hard.} \emph{Why:} the restricted-product topology (``almost all coordinates in $\mathcal{O}_v$'') is unlike any topology seen before this course. \emph{Mitigation:} build up from finite sub-products (``$S$-adèles'') as a directed limit before taking the full restricted product.
\textbf{Example questions}
\begin{enumerate}[noitemsep]
  \item[Q1 (foundational).] Show $\prod'(\Q_p,\Zp)$ together with $\mathbb R$ is locally compact. \emph{Rubric: basic opens are products of opens at finitely many places times $\prod \Zp$ elsewhere; local compactness follows from compactness of $\Zp$.}
  \item[Q2 (advanced).] Show $\Q$ embeds as a discrete subring of $\A_\Q$. \emph{Rubric: find a neighborhood of $0$ meeting $\Q$ only at $0$, using that a rational number bounded archimedean and integral everywhere finite must be an integer.}
\end{enumerate}

%---------------------------------------------------------------------
\subsection{Topology, Local Compactness, and Self-Duality of Adèles}
\textbf{Core concepts}
\begin{itemize}[noitemsep]
  \item Local compactness of $\A_K$
  \item Pontryagin duality background (via the measure-theory prerequisite)
  \item Self-duality $\A_K \cong \widehat{\A_K}$ via a nontrivial additive character
\end{itemize}
\textbf{Learning objectives}
\begin{itemize}[noitemsep]
  \item Prove local compactness of $\A_K$ from the restricted-product topology
  \item Explain the construction of self-duality via an additive character
  \item Apply Pontryagin duality formalism in the adelic setting
\end{itemize}
\textbf{Historically hard.} \emph{Why:} self-duality via a character is often presented too quickly, obscuring why the choice of character matters. \emph{Mitigation:} construct the standard character on $\A_\Q/\Q$ explicitly, local component by local component, before asserting self-duality in general.
\textbf{Example questions}
\begin{enumerate}[noitemsep]
  \item[Q1 (foundational).] Show $\A_K/K$ is compact, given an explicit fundamental domain. \emph{Rubric: exhibit a compact fundamental domain as a product of bounded local pieces.}
  \item[Q2 (advanced).] Construct $\psi = \prod \psi_v$ on $\A_\Q/\Q$ explicitly and verify it is nontrivial on $\A_\Q$ but trivial on $\Q$. \emph{Rubric: $\psi_\infty(x)=e^{-2\pi i x}$, $\psi_p(x)=e^{2\pi i \{x\}_p}$; sum of fractional parts over all places is an integer.}
\end{enumerate}

%---------------------------------------------------------------------
\subsection{The Idèle Group and Idèle Class Group}
\textbf{Core concepts}
\begin{itemize}[noitemsep]
  \item Idèles $\A_K^*$ as units of $\A_K$, with their own finer restricted-product topology
  \item The idèle class group $C_K = \A_K^*/K^*$
  \item Relation between $C_K$ and the classical ideal class group $\mathrm{Cl}(K)$
\end{itemize}
\textbf{Learning objectives}
\begin{itemize}[noitemsep]
  \item Explain why the idèle topology is strictly finer than the subspace topology from $\A_K$
  \item Derive the surjection from idèles to fractional ideals and identify its kernel
  \item Relate the idèle class group to the classical ideal class group
\end{itemize}
\textbf{Example questions}
\begin{enumerate}[noitemsep]
  \item[Q1 (foundational).] Show $(x_v) \mapsto \prod P_v^{v(x_v)}$ is a surjective homomorphism from idèles to fractional ideals, and identify the kernel. \emph{Answer: kernel is the group of ``unit idèles.''}
  \item[Q2 (advanced).] Show $C_K$ surjects onto $\mathrm{Cl}(K)$, and identify the kernel. \emph{Rubric: compose Q1's surjection with the quotient by principal ideals; kernel involves global units and archimedean components.}
\end{enumerate}

%---------------------------------------------------------------------
\subsection{Compactness of the Norm-One Idèle Class Group}
\textbf{Core concepts}
\begin{itemize}[noitemsep]
  \item The norm map $|\cdot|: \A_K^* \to \mathbb{R}_{>0}$; norm-one idèles $\A_K^1$
  \item Theorem: $C_K^1 = \A_K^1/K^*$ is compact
  \item This single theorem encodes both finiteness of the class number and Dirichlet's unit theorem
\end{itemize}
\textbf{Learning objectives}
\begin{itemize}[noitemsep]
  \item Prove compactness of $C_K^1$, building on Minkowski theory
  \item Explain how this theorem unifies finiteness of $\mathrm{Cl}(K)$ and the unit theorem
  \item Apply compactness to derive structural facts about $C_K$
\end{itemize}
\textbf{Historically hard.} \emph{Why:} this is the conceptual climax of the local-global story, and its proof repackages Minkowski theory and Dirichlet's unit theorem in idelic language --- students who haven't internalized those two results will find this proof opaque. \emph{Mitigation:} give an explicit dictionary mapping each step of the classical arguments to its idelic reformulation before the proof.
\textbf{Example questions}
\begin{enumerate}[noitemsep]
  \item[Q1 (foundational).] Explain why compactness of $C_K^1$ implies finiteness of the class number. \emph{Rubric: continuous image of compact is compact; the ideal-valued map has finitely many fiber classes.}
  \item[Q2 (advanced).] Explain how compactness of $C_K^1$ also encodes Dirichlet's unit theorem. \emph{Rubric: relate the kernel of the archimedean norm map on units to the logarithmic embedding lattice from \S2.6.}
\end{enumerate}

%---------------------------------------------------------------------
\subsection{Haar Measure, Fundamental Domains, and Volume Computations}
\textbf{Core concepts}
\begin{itemize}[noitemsep]
  \item Existence/uniqueness of Haar measure on $\A_K$ and $\A_K^*$
  \item An explicit fundamental domain for $K$ in $\A_K$; volume in terms of the discriminant
  \item Measure normalization conventions used later in Tate's thesis
\end{itemize}
\textbf{Learning objectives}
\begin{itemize}[noitemsep]
  \item Construct Haar measure on $\A_K$ from local Haar measures at each place
  \item Compute $\mathrm{vol}(K\backslash\A_K)$ in terms of the discriminant $d_K$
  \item Apply measure normalization conventions needed for zeta integrals
\end{itemize}
\textbf{Example questions}
\begin{enumerate}[noitemsep]
  \item[Q1 (foundational).] With $\mathrm{vol}(\Zp)=1$ at each finite place, compute $\mathrm{vol}(\A_\Q/\Q)$. \emph{Answer: volume $=1$, matching $d_\Q=1$ up to normalization.}
  \item[Q2 (advanced).] Express $\mathrm{vol}(\A_K/K)$ in terms of $d_K$ for general $K$, and explain its role in Tate's thesis. \emph{Rubric: $\mathrm{vol} = |d_K|^{1/2}$ up to convention; enters as a normalization factor in the residue of the global zeta integral at $s=1$.}
\end{enumerate}

%=====================================================================
\section{Module 4: Tate's Thesis (If Time Permits)}
\textbf{Weeks 13--15 (Sessions 25--30, contingent on pace).} \textit{Goal:} use the adelic/idelic machinery to give a uniform, harmonic-analytic proof of analytic continuation and the functional equation for the Dedekind zeta function and Hecke $L$-functions, following Tate's local-to-global method.

%---------------------------------------------------------------------
\subsection{Schwartz--Bruhat Functions and Local Zeta Integrals}
\textbf{Core concepts}
\begin{itemize}[noitemsep]
  \item Schwartz--Bruhat spaces at archimedean and non-archimedean places
  \item Local zeta integrals $Z(f,\chi,s) = \int f(x)\,\chi(x)\,|x|^s\, d^*x$
  \item Local convergence, meromorphic continuation, and local $L$-factors
\end{itemize}
\textbf{Learning objectives}
\begin{itemize}[noitemsep]
  \item Define Schwartz--Bruhat functions at each type of place with explicit examples
  \item Compute local zeta integrals for standard test functions
  \item Derive local $L$-factors from explicit computations at unramified places
\end{itemize}
\textbf{Example questions}
\begin{enumerate}[noitemsep]
  \item[Q1 (foundational).] Compute $Z(f,\chi,s)$ for $f = 1_{\mathcal{O}_p}$, $\chi$ unramified, non-archimedean. \emph{Answer: a geometric series yielding $(1-\chi(\pi)q^{-s})^{-1}$, the expected local $L$-factor.}
  \item[Q2 (advanced).] Compute the archimedean zeta integral for $f(x)=e^{-\pi x^2}$ with $\chi$ trivial. \emph{Rubric: Gaussian integral yielding a factor of $\Gamma(s/2)$ up to normalization.}
\end{enumerate}

%---------------------------------------------------------------------
\subsection{Local Functional Equation and Local $L$-Factors}
\textbf{Core concepts}
\begin{itemize}[noitemsep]
  \item Local Fourier transform and self-dual Haar measure
  \item The local functional equation relating $Z(f,\chi,s)$ and $Z(\hat f,\chi^{-1},1-s)$
  \item Local epsilon factors
\end{itemize}
\textbf{Learning objectives}
\begin{itemize}[noitemsep]
  \item State and explain the local functional equation for zeta integrals
  \item Derive the local functional equation at unramified places via explicit Fourier transforms
  \item Apply the local functional equation to compute epsilon factors in simple cases
\end{itemize}
\textbf{Historically hard.} \emph{Why:} the Fourier-analytic manipulation is analytically unfamiliar for students whose recent training has been algebraic, and epsilon factors are often introduced as an unmotivated ``fudge factor.'' \emph{Mitigation:} compute the functional equation fully and explicitly for $K=\Q_p$, unramified, before the general statement, showing exactly where the epsilon factor arises.
\textbf{Example questions}
\begin{enumerate}[noitemsep]
  \item[Q1 (foundational).] Verify the local functional equation at an unramified place directly from the \S4.1 computation. \emph{Rubric: compare the two local $L$-factors; epsilon factor equals $1$ in the unramified case.}
  \item[Q2 (advanced).] Compute the local epsilon factor for a ramified quadratic character of $\Q_p^*$. \emph{Rubric: a normalized Gauss-sum computation.}
\end{enumerate}

%---------------------------------------------------------------------
\subsection{Global Zeta Integrals, Poisson Summation, and the Global Functional Equation}
\textbf{Core concepts}
\begin{itemize}[noitemsep]
  \item The global zeta integral over idèle classes, factoring as a product of local integrals
  \item Poisson summation on $\A_K$, using self-duality (\S3.2) and the fundamental domain (\S3.5)
  \item Analytic continuation and functional equation; recovering $\zeta_K(s)$ and Hecke $L$-functions
\end{itemize}
\textbf{Learning objectives}
\begin{itemize}[noitemsep]
  \item Derive analytic continuation of the global zeta integral using Poisson summation
  \item Explain how the global functional equation follows from local functional equations plus Poisson summation
  \item Apply Tate's method to recover the classical functional equation of $\zeta_K(s)$
\end{itemize}
\textbf{Historically hard.} \emph{Why:} this step synthesizes nearly the entire course into one argument, and adelic Poisson summation is a non-obvious generalization of the classical formula on $\mathbb{R}$. \emph{Mitigation:} review classical Poisson summation on $\mathbb{R}$ as a warm-up, then present the adelic version as the same formula with $K$ replacing $\Z$ and $\A_K$ replacing $\mathbb{R}$.
\textbf{Example questions}
\begin{enumerate}[noitemsep]
  \item[Q1 (foundational).] State classical Poisson summation on $\mathbb{R}$ and give its adelic dictionary. \emph{Rubric: match each classical term (sum over $\Z$, Fourier transform on $\mathbb R$) to its adelic analogue.}
  \item[Q2 (advanced).] Sketch how Tate's method recovers the functional equation of $\zeta(s)$ for $K=\Q$. \emph{Rubric: split the zeta integral at $|x|=1$, apply adelic Poisson summation to relate the two pieces, obtain continuation and $\xi(s)=\xi(1-s)$.}
\end{enumerate}

%=====================================================================
\section{Prerequisite Map}
\begin{longtable}{p{3.6cm} p{6.2cm} p{4.4cm}}
\toprule
\textbf{Topic} & \textbf{Prerequisite(s)} & \textbf{Where taught} \\
\midrule
\endhead
1.1 Dedekind domains & Noetherian rings, localization, integral closure & Assumed background (algebra quals) \\
1.2 Fractional ideals / class group & Dedekind domain characterizations & \S1.1 \\
1.3 Extensions, decomposition & Unique factorization of ideals; integral closure & \S1.2; assumed background \\
1.4 Fundamental identity & Prime factorization in extensions & \S1.3 \\
1.5 Discriminant / different & Ramification index/residue degree; trace, norm & \S1.4; assumed background (Galois theory) \\
1.6 Decomposition/inertia groups & Discriminant and ramification; Galois theory & \S1.5; assumed background \\
2.1 Global fields & Ideal factorization and ramification framework & Module 1 \\
2.2 Absolute values / completions & Field theory; Cauchy sequences, completions & Assumed background (analysis quals) \\
2.3 Local field structure, Hensel & $p$-adic numbers; ramification concepts & \S2.2, \S1.4 \\
2.4 Local-global, product formula & Places of a global field; local field construction & \S2.1, \S2.2 \\
2.5 Minkowski theory & Number field structure; linear algebra / volumes & \S1.1, \S2.1; assumed background (algebra quals) \\
2.6 Dirichlet unit theorem & Minkowski theory / lattices; places & \S2.5, \S2.1 \\
3.1 Restricted direct products & Local fields; places of a global field & \S2.1--\S2.3 \\
3.2 Topology, self-duality & Restricted direct product topology; measure theory & \S3.1; assumed background (Haar measure) \\
3.3 Idèle / idèle class group & Adèle ring; ideal class group & \S3.1, \S1.2 \\
3.4 Compactness of $C_K^1$ & Idèle class group; Minkowski theory; unit theorem & \S3.3, \S2.5, \S2.6 \\
3.5 Haar measure, fundamental domain & Adèle/idèle topology; discriminant; measure theory & \S3.1, \S1.5; assumed background \\
4.1 Schwartz--Bruhat, local zeta integrals & Local fields; Haar measure on adèles/idèles & \S2.2, \S3.5; assumed background (real analysis) \\
4.2 Local functional equation & Local zeta integrals; self-duality (local analogue) & \S4.1, \S3.2 \\
4.3 Global functional equation & Local functional equation; Poisson summation prerequisites & \S4.2, \S3.2, \S3.5 \\
\bottomrule
\end{longtable}

%=====================================================================
\section{Assessment Plan}
Assessment philosophy: weekly homework only, no exams or final. Each set is matched to that week's module content and designed to catch specific failure modes, not merely to test recall.

\begin{longtable}{p{2.4cm} p{5.2cm} p{2.0cm} p{5.0cm}}
\toprule
\textbf{Module} & \textbf{Homework focus} & \textbf{Timing} & \textbf{Rationale} \\
\midrule
\endhead
Module 1 & Ideal factorization, ramification computations in explicit extensions, verifying the fundamental identity & Weeks 1--4 & Tests whether students can \emph{execute} the Dedekind-domain machinery computationally, not just recite definitions \\
Module 2 & $p$-adic computations, Hensel's Lemma applications, Minkowski bound and unit group computations & Weeks 5--9 & Tests the translation between local and global computations and geometry-of-numbers arguments, where reasoning errors typically surface \\
Module 3 & Adèle/idèle topology proofs, measure/volume computations, compactness argument write-ups & Weeks 10--12 & This material is proof-heavy and abstract; writing forces internalization of the topological/measure-theoretic arguments behind compactness \\
Module 4 & Local zeta integral computations; deriving the functional equation in worked cases (e.g.\ $K=\Q$) & Weeks 13--15 (contingent) & Since this module is only covered if time permits, homework is compressible and focuses on verifying the functional equation in the simplest nontrivial case as the essential takeaway \\
\bottomrule
\end{longtable}

\end{document}